\documentclass[11pt, a4paper]{article}

% ─── Packages ────────────────────────────────────────────────────────────────
\usepackage[T1]{fontenc}
\usepackage[utf8]{inputenc}
\usepackage[french]{babel}
\usepackage[top=2cm, bottom=2cm, left=2.5cm, right=2.5cm]{geometry}
\usepackage{booktabs}
\usepackage{xcolor}
\usepackage{titlesec}
\usepackage{amsmath}
\usepackage{hyperref}
\usepackage{array}
\usepackage{colortbl}
\usepackage{parskip}

% ─── Couleurs ────────────────────────────────────────────────────────────────
\definecolor{bleufonce}{RGB}{26, 55, 108}
\definecolor{bleuclair}{RGB}{39, 111, 191}
\definecolor{vert}{RGB}{46, 134, 90}
\definecolor{orange}{RGB}{224, 123, 34}
\definecolor{violet}{RGB}{106, 61, 154}
\definecolor{grisclair}{RGB}{242, 244, 248}
\definecolor{entete}{RGB}{26, 55, 108}

% ─── Style des titres ────────────────────────────────────────────────────────
\titleformat{\section}
  {\large\bfseries\color{bleufonce}}
  {\thesection.}{0.5em}{}
  [\vspace{-0.3em}\textcolor{bleufonce}{\rule{\linewidth}{0.8pt}}]

\titleformat{\subsection}
  {\normalsize\bfseries\color{bleuclair}}
  {\thesubsection}{0.5em}{}

\titlespacing{\section}{0pt}{14pt}{6pt}
\titlespacing{\subsection}{0pt}{8pt}{3pt}

% ─── Commandes tableaux ──────────────────────────────────────────────────────
\newcommand{\entetetab}[1]{\cellcolor{entete}\textcolor{white}{\textbf{#1}}}
\newcommand{\ligneA}[1]{\cellcolor{grisclair}{#1}}

% ─── Liens PDF ───────────────────────────────────────────────────────────────
\hypersetup{
  colorlinks=true,
  linkcolor=bleufonce,
  urlcolor=bleuclair,
  pdftitle={Rapport — Prévision des nappes phréatiques},
}

% ─────────────────────────────────────────────────────────────────────────────
\begin{document}

% ─── Titre ───────────────────────────────────────────────────────────────────
\begin{center}
  {\LARGE\bfseries\color{bleufonce}
    Prévision du niveau des nappes phréatiques\\[4pt]
    par réseaux de neurones}\\[8pt]
  {\normalsize Marius \quad — \quad Master 1 \quad — \quad 2025/2026}
\end{center}

\vspace{4pt}
\textcolor{bleufonce}{\rule{\linewidth}{1.2pt}}
\vspace{6pt}

% ═════════════════════════════════════════════════════════════════════════════
\section{Contexte et données}
% ═════════════════════════════════════════════════════════════════════════════

L'objectif est de prédire le niveau des nappes phréatiques (\textbf{GWL}, en mètres)
à partir de variables environnementales mesurées mensuellement sur des puits australiens.

\subsection{Jeu de données}

\begin{itemize}
  \setlength\itemsep{2pt}
  \item \textbf{2\,510 fichiers CSV} — un par puit — couvrant l'ensemble du territoire australien
  \item \textbf{8 régions :} VIC (1\,113 puits), NSW (687), WA (262), QLD (203),
        SA (198), TAS (19), NT (19), ACT (9)
  \item \textbf{6 variables par puit :} GWL (cible), Date, P (précipitations mm/mois),
        T (température °C), ET (évapotranspiration mm/mois), NDVI (indice de végétation)
  \item \textbf{OUVRAGES.csv :} identifiant, latitude, longitude et région de chaque puit
\end{itemize}

\textit{Note : de nombreux puits présentent des valeurs manquantes (NaN) dans la colonne GWL,
réduisant le nombre de puits exploitables selon les parties.}

% ═════════════════════════════════════════════════════════════════════════════
\section{Partie 1 — Visualisation}
% ═════════════════════════════════════════════════════════════════════════════

\subsection{Approche}

Une fonction générique \texttt{plot\_well(well\_id)} a été développée pour visualiser
n'importe quel puit. Elle génère deux figures de 5 sous-graphiques chacune (un par variable) :

\begin{itemize}
  \setlength\itemsep{2pt}
  \item \textbf{Figure A — Valeurs brutes :} courbes avec les unités réelles
        (m, mm/mois, °C, sans unité)
  \item \textbf{Figure B — Valeurs normalisées :} z-score appliqué colonne par colonne,
        permettant de comparer des variables d'unités différentes sur un même axe
\end{itemize}

La normalisation z-score est définie par :
\[
  \tilde{x} = \frac{x - \mu}{\sigma}
  \qquad \text{avec } \mu = \text{moyenne},\ \sigma = \text{écart-type de la séquence}
\]

La fonction \texttt{plot\_sample()} sélectionne automatiquement $n$ puits par région
via \texttt{OUVRAGES.csv} et sauvegarde les graphiques au format PNG.

% ═════════════════════════════════════════════════════════════════════════════
\section{Partie 2a — Modèle LSTM générique}
% ═════════════════════════════════════════════════════════════════════════════

\subsection{Approche}

Un unique réseau LSTM est entraîné sur l'ensemble des puits australiens (modèle générique).
Il apprend à prédire GWL à partir des variables environnementales passées.

\begin{itemize}
  \setlength\itemsep{2pt}
  \item \textbf{Découpage train/test :} toute la chronique moins les 12 derniers mois
        pour l'entraînement, les 12 derniers mois pour le test (conforme aux consignes)
  \item \textbf{Features d'entrée :} P, T, ET, NDVI + latitude normalisée + longitude
        normalisée (6 variables)
  \item \textbf{Cible :} GWL au mois suivant
  \item \textbf{Séquences glissantes} de 12 mois (\texttt{SEQ\_LEN = 12})
  \item \textbf{Normalisation z-score} par puit et par variable, calculée sur les données
        d'entraînement uniquement
\end{itemize}

\subsection{Architecture}

\begin{itemize}
  \setlength\itemsep{2pt}
  \item Type : \textbf{LSTM} (Long Short-Term Memory) — réseau récurrent avec mémoire
        à long terme
  \item Configuration finale (issue du grid search) : 3 couches, hidden = 128,
        dropout = 0.2
  \item Couche de sortie : fully connected ($128 \to 1$) — prédit une valeur scalaire GWL
  \item Optimiseur : Adam \quad|\quad Fonction de perte : MSE
\end{itemize}

\subsection{Grid Search — 108 combinaisons testées}

Recherche exhaustive sur la grille suivante, résultats sauvegardés dans
\texttt{grid\_search\_results.csv} :

\vspace{6pt}
\begin{center}
\begin{tabular}{>{\bfseries}lll}
  \toprule
  \entetetab{Hyperparamètre} & \entetetab{Valeurs testées} & \entetetab{Meilleure valeur} \\
  \midrule
  \ligneA{hidden (neurones/couche)} & \ligneA{64, 128, 256}       & \ligneA{128}   \\
  n\_layers (couches LSTM)          & 1, 2, 3                     & 3              \\
  \ligneA{lr (learning rate)}       & \ligneA{1e-3, 5e-4, 1e-4}  & \ligneA{5e-4}  \\
  batch\_size                        & 256, 512                    & 512            \\
  \ligneA{epochs}                   & \ligneA{30, 50}             & \ligneA{30}    \\
  \bottomrule
\end{tabular}
\end{center}
\vspace{6pt}

\subsection{Optimisation GPU}

Le dataset complet a été préchargé en VRAM (RTX 5070 Ti — 16 Go) sous forme de tenseurs
PyTorch. Le batching est effectué directement sur la GPU via \texttt{torch.randperm()},
éliminant tout transfert CPU $\to$ GPU pendant l'entraînement.

\subsection{Résultats — 854 puits}

\vspace{4pt}
\begin{center}
\begin{tabular}{lrrr}
  \toprule
  \entetetab{Métrique} & \entetetab{Moyenne} & \entetetab{Écart-type} & \entetetab{Médiane} \\
  \midrule
  \ligneA{RMSE (m)} & \ligneA{0.805} & \ligneA{$\pm$0.968} & \ligneA{0.440} \\
  MAE (m)           & 0.686          & $\pm$0.869          & 0.356          \\
  \bottomrule
\end{tabular}
\end{center}
\vspace{4pt}

\textit{La médiane est privilégiée : l'écart-type élevé indique la présence de quelques puits
très difficiles à prédire qui tirent la moyenne vers le haut.}

% ═════════════════════════════════════════════════════════════════════════════
\section{Partie 2b — Chronos 2 (Foundation Model)}
% ═════════════════════════════════════════════════════════════════════════════

\subsection{Approche}

Utilisation du modèle pré-entraîné \texttt{amazon/chronos-bolt-small} issu de la librairie
\texttt{chronos-forecasting} (v2.2.2). Aucun entraînement sur nos données.

\begin{itemize}
  \setlength\itemsep{2pt}
  \item \textbf{Modèle univarié :} seul l'historique GWL est utilisé comme contexte
        (pas P, T, ET, NDVI)
  \item \textbf{Prédiction zero-shot :} le modèle prédit directement sans fine-tuning
  \item \textbf{Contexte :} toutes les valeurs GWL non-NaN avant les 12 derniers mois
  \item \textbf{Prédiction probabiliste :} $N$ scénarios générés $\to$ médiane retenue
        comme prédiction finale
  \item Traitement par batch de 32 puits en parallèle sur GPU
  \item Critère d'inclusion : minimum 24 mois de GWL non-NaN dans la période d'entraînement
\end{itemize}

\subsection{Résultats — 1\,174 puits}

\vspace{4pt}
\begin{center}
\begin{tabular}{lrrr}
  \toprule
  \entetetab{Métrique} & \entetetab{Moyenne} & \entetetab{Écart-type} & \entetetab{Médiane} \\
  \midrule
  \ligneA{RMSE (m)} & \ligneA{0.673} & \ligneA{$\pm$0.973} & \ligneA{0.301} \\
  MAE (m)           & 0.578          & $\pm$0.847          & 0.248          \\
  \bottomrule
\end{tabular}
\end{center}

% ═════════════════════════════════════════════════════════════════════════════
\section{Comparaison et conclusion}
% ═════════════════════════════════════════════════════════════════════════════

\vspace{4pt}
\begin{center}
\begin{tabular}{lrrr}
  \toprule
  \entetetab{Métrique} & \entetetab{LSTM (2a)} & \entetetab{Chronos 2 (2b)}
    & \entetetab{Gain} \\
  \midrule
  \ligneA{RMSE moyen (m)}  & \ligneA{0.805} & \ligneA{0.673}   & \ligneA{\textcolor{vert}{$-16\,\%$}} \\
  MAE moyen (m)            & 0.686          & 0.578            & \textcolor{vert}{$-16\,\%$}          \\
  \ligneA{RMSE médian (m)} & \ligneA{0.440} & \ligneA{0.301}   & \ligneA{\textcolor{vert}{$-32\,\%$}} \\
  MAE médian (m)           & 0.356          & 0.248            & \textcolor{vert}{$-30\,\%$}          \\
  \ligneA{Puits évalués}   & \ligneA{854}   & \ligneA{1\,174}  & \ligneA{—}                           \\
  \bottomrule
\end{tabular}
\end{center}
\vspace{6pt}

Chronos 2 surpasse le LSTM sur l'ensemble des métriques, avec une réduction d'erreur
de 30 à 32\,\% sur les médianes. Ce résultat est d'autant plus remarquable que :

\begin{itemize}
  \setlength\itemsep{2pt}
  \item Chronos n'utilise que l'historique GWL (univarié), sans accès aux variables
        météo ni GPS
  \item Aucun entraînement sur les données australiennes n'a été effectué (zero-shot)
  \item Le LSTM dispose pourtant de 6 features d'entrée et de 30 epochs d'entraînement
\end{itemize}

\textbf{Explication :} les nappes phréatiques présentent des patterns temporels réguliers
(cycles saisonniers, tendances lentes) que Chronos 2 reconnaît grâce à son pré-entraînement
sur des milliers de séries temporelles diverses. Le LSTM doit apprendre ces patterns depuis
zéro avec un volume de données relatif limité.

\textit{Limite de la comparaison :} les ensembles de puits évalués diffèrent (854 pour le
LSTM, 1\,174 pour Chronos) en raison de critères d'inclusion différents. Une comparaison
stricte sur le même sous-ensemble renforcerait la validité des conclusions.

\end{document}
